\section{Giới thiệu}

Dự án này khai thác bộ dữ liệu World Development Indicators (WDI) của World Bank để phân tích mức độ phát triển của các quốc gia và khu vực trong giai đoạn 2000--2023. Mục tiêu tổng quan là xây dựng một bức tranh trực quan về phát triển bền vững, trong đó tăng trưởng kinh tế được đặt cạnh các kết quả về giáo dục, sức khỏe--dân số và môi trường. Cách tiếp cận này giúp tránh việc đánh giá phát triển chỉ bằng một chỉ số đơn lẻ, đồng thời làm rõ sự khác biệt giữa các khu vực, nhóm thu nhập và quốc gia theo thời gian.

Các phân tích được triển khai trên Tableau sau khi dữ liệu được tiền xử lý từ nguồn WDI gốc. Bốn hướng phân tích gồm phát triển kinh tế, giáo dục, sức khỏe--dân số, và môi trường--biến đổi khí hậu được dùng như các góc nhìn bổ trợ để trả lời cùng một bài toán chung. Trong báo cáo, kết quả trực quan hóa được diễn giải theo hướng mô tả khách quan, phân biệt rõ giữa mẫu hình quan sát được, nhận định phân tích và giới hạn của dữ liệu.

\subsection{Thành viên và phân công}

\begin{longtable}{L{3cm}L{5cm}L{6cm}}
\label{tab:members}\\
\toprule
MSSV & Họ tên & Phần phụ trách \\
\midrule
\endfirsthead
\toprule
MSSV & Họ tên & Phần phụ trách \\
\midrule
\endhead
23120237 & Lê Lâm Trí Đức & Tổng quan dataset, phân tích phát triển kinh tế, bố cục dashboard tổng \\
23120189 & Hoàng Quốc Việt & Thống kê mô tả, phân tích giáo dục, định dạng màu sắc \\
23120190 & Nguyễn Lê Thế Vinh & Phân tích sức khỏe và dân số, filter/selector, video giới thiệu Tableau \\
23120202 & Phạm Quang Vinh & Phân tích môi trường, dashboard actions, video trình bày phân tích \\
\bottomrule
\end{longtable}

\subsection{Mục tiêu báo cáo}

Báo cáo hướng đến ba mục tiêu chính. Thứ nhất, các tập dữ liệu sạch được xây dựng từ WDI cho giai đoạn 2000--2023 theo cùng một quy ước tiền xử lý. Thứ hai, các tập dữ liệu này được trực quan hóa trong Tableau để so sánh giữa quốc gia, khu vực, nhóm thu nhập và thời gian. Thứ ba, kết quả trực quan được diễn giải trong một khung phân tích chung về phát triển bền vững, trong đó mỗi mục tiêu thành phần cung cấp một nhóm bằng chứng riêng nhưng cùng phục vụ câu hỏi nghiên cứu trung tâm. Cách tổ chức này tạo được mạch lập luận thống nhất từ dữ liệu, phương pháp, kết quả đến thảo luận.
